\section{Uvod}

U ovom seminarskom radu rješavaju se četiri zadatka iz kolegija \textit{Optimizacijske metode u razvoju i upravljanju sustava}. Redom su to: kvadratni program s dva ograničenja nejednakosti, njegova robusna inačica uz intervalnu nesigurnost koeficijenata, formulacija nivelacije uzdužnog profila ceste kao linearnog programa te sinteza regulatora po varijablama stanja za sklopni linearni sustav pomoću linearnih matričnih nejednakosti.

Zajednička nit svih zadataka je konveksnost: nakon što se problem zapiše u standardnom obliku (LP, QP ili SDP), lokalni optimum je ujedno globalan, a KKT uvjeti odnosno linearne matrične nejednakosti daju provjerljiv certifikat optimalnosti ili stabilnosti~\cite{boyd2004convex,jokic2024predavanja}.

Svi proračuni provedeni su u MATLAB-u (verzija R2025b): \texttt{quadprog} za kvadratne programe u zadacima 1 i 2, \texttt{linprog} za linearni program u zadatku 3 te YALMIP sa rješavačem SDPT3 za semidefinitni program u zadatku 4. Skripte se nalaze u prilogu; svaka brojka i svaka slika u radu proizlazi iz jednog pokretanja tih skripti.

Poglavlja slijede numeraciju zadataka:
\begin{itemize}
    \item \textbf{Poglavlje 2} --- Zadatak 1: konveksnost, KKT uvjeti, \texttt{quadprog}, Lagrangeov dualni problem i osjetljivost optimalne vrijednosti.
    \item \textbf{Poglavlje 3} --- Zadatak 2: robusni ekvivalent uz nesigurnost koeficijenata od $\pm 15\%$ i usporedba s nominalnim rješenjem.
    \item \textbf{Poglavlje 4} --- Zadatak 3: diskretizacija konačnim razlikama, LP formulacija i rješenja za četiri zadana slučaja.
    \item \textbf{Poglavlje 5} --- Zadatak 4: stabilnost pojedinih modova, simulacija sklopnog sustava i LMI sinteza stabilizirajućeg regulatora.
\end{itemize}
